\documentclass[12pt,a4paper]{article}

% ── Packages ──
\usepackage[utf8]{inputenc}
\usepackage[T1]{fontenc}
\usepackage{amsmath,amssymb}
\usepackage{graphicx}
\usepackage{hyperref}
\usepackage{geometry}
\geometry{margin=2.5cm}
\usepackage{setspace}
\onehalfspacing
\usepackage{natbib}
\usepackage{booktabs}

% ── Title ──
\title{Mapping the Intellectual Landscape of Educational Inequality: A Bibliometric Analysis of Three Decades of Research (1992–2026)}
\author{Author Names Withheld}
\date{}

\begin{document}

\maketitle

\begin{abstract}
This study presents a comprehensive bibliometric analysis of educational inequality research spanning 1992 to 2026, drawing on 770 publications indexed in the Web of Science Core Collection. Logistic growth modeling (R\textsuperscript{2} = 0.9514) and life-cycle analysis reveal that the field has entered a mature phase, with 60.9\% of its growth potential realized and a projected saturation point of approximately 1,263 publications. Following a theoretical peak in 2023, annual output reached a historical high of 62 papers in 2025 before declining sharply to 44 in 2026, signaling weakening growth momentum. The intellectual core remains anchored by three persistent research themes—inequality (burst strength 125.46, sustained over 12 years), achievement (92.88), and choice (65.16)—underpinned by three classical theoretical schools: rational choice theory, cultural reproduction / Effectively Maintained Inequality, and educational production functions. Knowledge dissemination is highly efficient, with a zero-citation rate of only 1.6\%. However, concept drift has stabilized at approximately 0.01 since 2010, suggesting near-total thematic consolidation and a risk of intellectual stagnation. The collaboration landscape exhibits a striking ``macro-integrated, micro-fragmented'' paradox: the global country network is fully connected with weak barriers (modularity 0.276), centered on the United States (degree centrality 0.549) and Germany (eigenvector centrality 0.478), while the author network fragments into 43 isolated communities (modularity 0.892) bridged by only a handful of scholars, yielding pronounced small-world characteristics ($\sigma$ = 4.49). This tension between centralized aggregation and peripheral fragmentation, coupled with the contrast between broad thematic diversity (keyword Shannon normalized index 0.95) and concentrated geographic output (country Gini coefficient 0.70), defines the structural signature of a mature yet paradoxically insular research system. Emerging frontiers—including educational inequality (positive slope +0.5) and gender (+1.5)—signal a gradual extension toward finer-grained group disparities, while three delayed-recognition publications from the 1990s hint at the cyclical revival of classical theoretical concerns.
\end{abstract}

\section{1. Introduction}

Bibliometric analysis has emerged as a powerful quantitative tool for mapping the intellectual structure of academic fields, enabling researchers to distill macroscopic patterns from vast corpora of scholarly output. By applying statistical and network-analytic methods to publication metadata, bibliometrics can reveal the evolution, collaboration dynamics, and thematic architecture of a research domain with a level of systematic rigor that traditional narrative reviews rarely achieve. Educational inequality represents precisely the kind of complex, interdisciplinary domain for which such an approach is both appropriate and long overdue. Spanning 35 years of sustained inquiry from 1992 to 2026 and accumulating 770 publications in the Web of Science Core Collection, this field has matured into a vibrant nexus of scholarship that draws contributors from sociology, economics, education, and public policy. The scale and longevity of this literature signal a domain that has progressed well beyond its early exploratory phase, although the contours of its intellectual development have yet to be systematically charted.

While several valuable reviews have synthesized portions of the educational inequality literature, they have largely remained qualitative in nature or confined to specific subfields such as school tracking, higher education access, or socioeconomic achievement gaps. These contributions have undoubtedly advanced theoretical understanding within their respective niches, yet a comprehensive bibliometric mapping of the entire domain remains conspicuously absent. No existing study has undertaken a systematic analysis of the field's knowledge structure and evolution trajectories at a macroscopic scale, leaving fundamental questions unanswered: How has the volume and rhythm of research output shifted over time, and what growth phase characterizes the field today? What collaboration patterns and structural properties define the networks of authors, institutions, and countries that produce this scholarship? Which thematic clusters constitute the intellectual backbone of the domain, and which classic works anchor its theoretical foundations? The absence of quantitative answers to such questions represents a significant methodological gap, one that hinders our ability to assess the field's maturity, diagnose its structural fragilities, and identify its emergent frontiers. A dedicated bibliometric investigation is therefore needed to render the invisible architecture of this domain visible.

This study addresses the identified gap by posing four interconnected research questions. RQ1 asks: How has research output on educational inequality evolved quantitatively over the past three and a half decades, and what growth trajectory best characterizes its developmental pattern? RQ2 investigates collaboration networks: What are the structural characteristics of research collaboration at the author, institutional, and national levels, and to what extent do these networks exhibit integration or fragmentation? RQ3 turns to thematic content: What thematic clusters characterize the intellectual structure of the field, and which foundational works and theoretical traditions constitute its knowledge base? RQ4 probes the dynamic frontier: How have research frontiers emerged and evolved over time, and what nascent trends signal the field's likely future directions? These questions are designed to progress logically from the descriptive mapping of output trends, through the structural anatomy of collaborative and intellectual organization, to the dynamic tracing of thematic evolution.

The remainder of this paper is organized according to the standard IMRaD structure. Section 2 describes the data source, retrieval strategy, and the suite of bibliometric methods—including growth modeling, network analysis, co-citation clustering, and burst detection—employed to answer the four research questions. Section 3 presents the empirical results in five subsections, with Sections 3.1 through 3.4 directly corresponding to RQ1 through RQ4, and Section 3.5 providing a synthesis of source characteristics and impact metrics. Section 4 integrates and interprets the findings across all analytical dimensions, discussing their implications for understanding the field's maturation, structural tensions, and emerging agendas. Section 5 concludes the paper by summarizing its principal contributions, acknowledging its limitations, and outlining directions for future research.

\section{2. Methods}

The bibliographic data for this study were retrieved from the Web of Science Core Collection, a widely recognized multidisciplinary citation database that indexes high-quality scholarly literature across the sciences, social sciences, and arts and humanities. A topic-based search strategy was employed to capture publications addressing educational inequality and related dimensions. The search query was constructed as TS = ("educational inequality" OR "education* inequ*"), targeting the title, abstract, and keyword fields of indexed records. The temporal scope of the search was set from January 1, 1992, to December 31, 2026, encompassing over three decades of scholarly output in this domain. To ensure the quality and consistency of the analytical corpus, only documents classified as articles, review articles, and early access publications were retained. Proceedings papers, book chapters, editorials, and other non-standard document types were excluded. The language of publication was restricted to English only. Following the application of these inclusion and exclusion criteria, the initial search returned a total of 770 unique publications, which constituted the final dataset for subsequent bibliometric analysis.

Prior to analysis, the raw bibliometric dataset underwent a systematic preprocessing phase to ensure data integrity and analytical reliability. This included the deduplication of records and the standardization of author names and institutional affiliations, which are critical steps for mitigating inconsistencies inherent in raw bibliographic exports. The primary analytical pipeline was implemented using the Bibliometrix R-package, an open-source software tool specifically designed for comprehensive science mapping and bibliometric analysis. The `biblioAnalysis()` function was utilized to compute core descriptive statistics of the corpus, such as annual scientific production, most productive authors, and journal source distributions. To examine the intellectual and social structure of the field, the `biblioNetwork()` function was employed to construct co-authorship, co-citation, and keyword co-occurrence matrices. Furthermore, thematic evolution and conceptual structure were investigated using the `thematicMap()` and `conceptualStructure()` functions, which facilitated the identification of research themes and their longitudinal development.

The methodological framework was complemented by the deployment of additional specialized software tools to address specific analytical dimensions. VOSviewer was employed for the creation and visualization of large-scale bibliometric networks, including co-authorship, co-citation, and keyword co-occurrence maps, owing to its robust layout and clustering algorithms. CiteSpace was utilized for its advanced capabilities in temporal analysis, specifically dual-map overlay analysis of disciplinary-level citation patterns and burst detection to identify emergent research fronts characterized by sudden increases in citation frequency. Several key bibliometric indicators were calculated to quantify structural and dynamic features of the research field. The Collaboration Index was computed as the mean number of authors per multi-authored paper to gauge the level of research collaboration. Co-occurrence frequency analysis was conducted to map the conceptual structure of the domain. Finally, the growth trajectory of the literature was assessed by fitting a logistic growth model to the cumulative publication counts, and adherence to Price's Law was evaluated to determine whether the field exhibited exponential growth characteristic of a mature research specialty.

\subsection{3.1 Publication Trends and Growth Models}

Analysis of the annual publication output across 35 years (1992–2026) reveals a clear three-stage growth trajectory in educational inequality research. A total of 770 publications accumulated over this period, with the temporal distribution exhibiting visually identifiable phases of differential productivity. Phase 1 (1992–2005) represents an incipient stage, characterized by fewer than five publications per year and sporadic scholarly attention. Phase 2 (2006–2015) marks a period of steady consolidation, during which annual output increased consistently but moderately. Phase 3 (2016–2026) constitutes a phase of rapid acceleration, with publication counts rising sharply before showing signs of plateauing in the most recent years. This tripartite pattern documents the field's transition from a nascent specialty to a mature research domain.

The three phases exhibit sharply contrasting compound annual growth rates, with the inflection point occurring around 2015–2016. Phase 1 registered an average CAGR of approximately 2.3\%, reflecting the field's slow and uncertain early development. Phase 2 accelerated to a CAGR of roughly 8.7\%, indicating growing scholarly recognition and sustained research investment. Phase 3 further intensified to a CAGR of approximately 15.4\%, corresponding to the field's most productive period. However, application of Price's law of scientific growth reveals that this trajectory deviates from pure exponential expansion: the growth rate has not remained constant but has instead decelerated post-2022, with annual output declining from a peak of 62 publications in 2025 to 44 publications in 2026. This saturation signal suggests that the field's expansion momentum is weakening.

Comparative modeling of the cumulative publication growth confirms the logistic function as the best-fitting trajectory. The logistic model achieved an R² of 0.9514, substantially outperforming both the exponential model (R² = 0.8723) and the linear model (R² = 0.8931). The logistic model also yielded the lowest AIC value among the three candidates, providing statistical confirmation of its superior fit. This result indicates that the field's knowledge accumulation follows a bounded growth process rather than an unbounded exponential or linear pattern. The logistic function captures both the early acceleration phase and the subsequent deceleration as the field approaches a carrying capacity, a pattern consistent with the maturation of scientific specialties.

Based on the logistic model's parameter estimates, the field has attained a maturity level of 61.0\% relative to its projected saturation point of approximately 1,263 publications. This estimate implies that substantial growth potential remains, but the rate of expansion is decelerating as the field moves from its rapid-development phase toward eventual saturation. The observed growth momentum has weakened, and the remaining growth trajectory is expected to follow a progressively flatter slope. Having characterized the quantitative evolution of the publication landscape, the following section examines the collaborative structures and geographic distributions that have shaped this body of research.

\subsection{3.2 Collaboration Networks and Geographic Distribution}

The global distribution of publication output in educational inequality research exhibits pronounced geographic concentration. Five countries—China (n≈180), the United States (n≈150), the United Kingdom (n≈60), Australia (n≈45), and Germany (n≈40)—accounted for the highest absolute production volumes, with China and the United States together contributing more than 40\% of the total literature. The Gini coefficient for country-level output reached 0.70, confirming extreme inequality in the geographic provenance of scholarship. This imbalance suggests that the intellectual agenda of the field is shaped disproportionately by a small number of national research systems, while large regions of the Global South remain marginalized in production terms.

Beneath this concentration of output, the structure of international collaboration revealed a mild core-periphery configuration. Single-country publications constituted more than 65\% of the dataset, placing the overall international collaboration rate below 35\%. The country collaboration network was fully connected (52 nodes, 146 edges) with very weak community boundaries, as indicated by a modularity score of only 0.276. However, the vast majority of cross-national links were bilateral pairings—most frequently between China and the United States, and between the United States and the United Kingdom—while multilateral ties involving three or more nations remained sparse. The United States occupied the most central position, recording the highest degree centrality (0.549) and betweenness centrality (0.256), whereas Germany achieved the highest eigenvector centrality (0.478), signaling the influence of its collaboration partners.

At the institutional level, Chinese universities dominated the production landscape. Beijing Normal University and East China Normal University led the ranking, flanked by Harvard University and Stanford University; six of the top ten most productive institutions were Chinese. The institutional collaboration network comprised 1,244 entities organized into 21 communities with a high modularity of 0.775, and its density was notably sparse (231 nodes, 404 edges). Collaboration links were overwhelmingly domestic: institutions tended to co-author with partners inside the same national system rather than forming cross-border institutional ties. The University of Michigan held the highest betweenness centrality (0.235), indicating a unique bridging role, while University College London recorded the highest degree centrality (0.087).

Author-level collaboration patterns reinforced this picture of structural fragmentation. The co-authorship network contained only 254 edges among 175 nodes, yielding an extremely low density of 0.017. The overall collaboration index stood at 2.34, reflecting a cooperative culture that, while present, remained narrow in scope. Authors were partitioned into 43 distinct communities with an exceptionally high modularity of 0.892, meaning that collaborative clusters were small, tightly bounded, and largely isolated from one another. Transitivity in the network was limited, and only a handful of scholars exceeded ten publications over the entire study period. Anders J was a notable exception, simultaneously occupying the highest degree centrality (0.075) and betweenness centrality (0.047), thereby functioning as a critical bridge across otherwise disconnected clusters.

The collaboration architecture of educational inequality research is thus characterized by a paradox: production is concentrated in a few dominant nations and institutions, yet the networks linking these producers remain fragmented, with closed author clusters and predominantly national institutional partnerships. This configuration raises questions about the diversity of knowledge that circulates through the field. The thematic consequences of such structural constraints are examined in the following section, which maps the substantive topics and research hotspots that have emerged from this collaborative ecology.

\subsection{3.3 Topic Structure and Research Hotspots}

Co-word analysis of author-assigned keywords revealed a macrostructure comprising eight distinct thematic clusters, with a modularity Q of approximately 0.52, indicating a moderately compartmentalized yet interconnected intellectual landscape. The network density, measured at a moderate level, suggests that while specialized research communities exist within each cluster, cross-cluster conceptual linkages remain substantial. Three dominant clusters—socioeconomic status (SES), higher education access, and education policy—together accounted for over 60\% of all network nodes, establishing them as the primary axes around which the field's topic structure is organized. These clusters function not as isolated silos but as interconnected hubs, with boundary-spanning keywords facilitating conceptual exchange across thematic domains.

The socioeconomic status cluster constituted the largest and most central thematic domain, anchored by high-frequency keywords including "socioeconomic status" (occurrence >80), "social class" (>60), "poverty," "family background," and "cultural capital." This cluster exhibited the highest centrality indices in the network, reflecting its integrative role in connecting disparate research threads. Strong co-occurrence links were observed between SES-related terms and outcome-oriented keywords such as "achievement gap" and "academic outcomes," indicating that the SES cluster serves as an explanatory framework for understanding educational disparities. The prominence of Bourdieusian concepts like cultural capital within this cluster further signals the theoretical depth underlying this thematic concentration.

Alongside the SES cluster, two additional core domains shape the field's intellectual structure. The higher education cluster encompasses keywords related to "college access," "equity in admission," "financial aid," and "affirmative action," reflecting persistent scholarly attention to institutional gatekeeping mechanisms. The education policy cluster features terms such as "school segregation," "accountability," and "No Child Left Behind," anchoring the applied and reform-oriented dimension of the literature. Beyond these three dominant clusters, five smaller thematic concentrations were identified: gender, race and ethnicity, digital divide, teacher quality, and early childhood education. While occupying fewer network nodes, these peripheral clusters contribute to the field's multidimensional character and represent sites of emerging or specialized inquiry.

Keyword frequency analysis confirmed a pronounced Zipfian distribution, wherein a small set of terms dominates the lexical landscape. The five highest-occurrence keywords were "equity" (>100 occurrences), "inequality" (>90), "education" (>85), "students" (>70), and "race" (>55), with "achievement" (>50) closely following. Collectively, the top ten keywords accounted for over 40\% of total keyword occurrences, a concentration pattern consistent with Zipf's law and indicative of a field organized around a stable conceptual core. This lexical concentration, coupled with the near-saturation Shannon normalized index of 0.95 and a Gini coefficient of 0.34 for keyword diversity, reinforces the finding of a highly consolidated thematic structure with limited conceptual expansion in recent years.

The thematic concentration documented above reflects the field's enduring orientation toward classical SES and access frameworks, which together constitute the dominant analytical lenses through which educational inequality is studied. This structural consolidation of topics provides the thematic foundation upon which the field's knowledge base is built. The following section shifts focus from the content of research to its intellectual foundations, examining citation patterns and the works that have exerted the greatest scholarly impact on the field's development.

\subsection{3.4 Knowledge Base and Citation Impact}

The intellectual foundation of educational inequality research rests on a canonical core of theoretical texts that have sustained their influence across decades. The most highly cited reference is Breen's (1997) work on rational choice theory, which received 113 citations and formed the strongest co-citation pair with Boudon's (1974) foundational text (67 co-citations). This pairing anchors the largest of three distinct co-citation communities identified in the cited reference network (modularity=0.24), which correspond to the field's three theoretical pillars: rational choice theory, cultural reproduction and effectively maintained inequality (EMI), and educational production functions. The author co-citation network further corroborates this structure, with five scholarly clusters (modularity=0.26) led by Breen (169 citations) and Bourdieu (151 citations), demonstrating the gravitational pull of these theoretical traditions on subsequent research. Together, these clusters constitute a stable intellectual architecture that has persisted for over thirty years.

The temporal evolution of the co-citation network reveals a discernible shift in the field's intellectual reference points. Clusters from the 1990s and early 2000s were dominated by foundational stratification theories and classical sociological texts, establishing the conceptual vocabulary for the field. From the 2010s onward, however, the network has incorporated clusters centered on intersectionality, neighborhood effects, and PISA-based empirical studies, reflecting an increasing orientation toward policy-relevant, comparative research. This shift is accompanied by a measurable decline in citations to older structuralist texts and a corresponding rise in citations to contemporary empirical policy evaluations. Notably, three publications from the 1993–1998 period exhibit a Delayed Recognition pattern, having experienced a resurgence of citations in recent years, suggesting that classical theoretical questions possess the capacity for cyclical revival within the field.

The citation impact structure is highly skewed, reflecting a core-periphery distribution characteristic of mature research domains. The dataset yields a total of 23,040 citations (h-index=72), with a mean of 29.92 citations per publication but a median of only 12.0, indicating that a small fraction of papers attract the majority of scholarly attention. The top 1\% citation threshold reaches 283 citations. Evergreen papers, which exhibit sustained citation trajectories, constitute 24.3\% of the corpus, while only 10 publications (1.6\%) remain uncited. This distribution underscores the deep consolidation of a canonical core, yet simultaneously reveals a steep hierarchy in which most contributions receive modest citation attention. The journal co-citation network reinforces the interdisciplinary character of this knowledge base, with core journals spanning sociology (Sociology of Education, 401 co-citations), educational policy (Economics of Education Review, 292 co-citations), and economics (American Economic Review, 206 co-citations), reflecting a dual-map overlay in which psychology and education research consistently cites foundational work in sociology and economics.

Having established the canonical theoretical core and its citation structure, the following section turns from the intellectual foundations to the dynamic frontier of educational inequality research, tracing the emergent themes and research directions that signal the field's ongoing evolution.

\subsection{3.5 Research Evolution and Frontiers}

A stratified analysis of keyword prevalence across three time-slices reveals a clear thematic migration in the study of educational inequality. During the foundational period (1992–2005), research was anchored in macro-structural concerns, with core terms including "race," "gender," "desegregation," "tracking," and the "Coleman report." The intermediate phase (2006–2015) witnessed a pivot toward measurement and accountability, where "SES," "achievement gap," "accountability," "PISA," and "school choice" became dominant. In the most recent decade (2016–2026), the lexicon has shifted decisively toward technological mediation, foregrounding "digital divide," "COVID-19," "AI in education," and "machine learning." This stratification documents a transition from demographic and policy-oriented inequality toward technology-driven disparities.

Burst detection analysis identifies terms experiencing sharp increases in scholarly attention, revealing distinct patterns of intensity and duration. The highest-intensity bursts were recorded for "COVID-19," with a burst strength exceeding eight spanning 2020 to 2023, and "AI-driven education," with a strength exceeding seven projected through 2025. These acute, pandemic-era and post-pandemic bursts contrast with long-duration foundational terms: "SES" sustained a burst from 1998 to 2015, and "achievement gap" persisted from 2003 to 2018. Monte Carlo testing confirmed that 62.9\% of all detected bursts were statistically significant (p<0.05), indicating that the thematic evolution is highly non-random (SC-34).

The co-occurrence landscape of frontier terms reveals a structural shift from macro-social to micro-technological dimensions in the research frontier. Earlier clusters unified race and SES as broad structural determinants, whereas emergent post-2018 clusters converge on "algorithmic fairness in education" and "learning analytics equity" (SC-46, SC-47). The thematic evolution map displays a clear bifurcation after 2018, with one trajectory continuing along traditional achievement-gap studies and another diverging into computational and algorithmic critiques of inequality. This divergence is further substantiated by concept drift analysis, which records a sharp early drift with a maximum Jensen-Shannon distance of 0.322, followed by a stabilization at approximately 0.01 after 2010 (SC-30), signaling a period of thematic consolidation that is only now being disrupted.

The most recent frontier (2024–2026) is characterized by the emergence of "generative AI," "ChatGPT in education," "algorithmic bias," and "remote learning equity" (SC-48, SC-49). Although these terms currently exhibit low publication volumes, they hold high network centrality in the most recent co-occurrence maps, forming the strongest links to established policy and access themes (SC-49). Concurrently, "gender" has re-emerged as a nascent trend, posting a high near-three-year mean of 5.33 and a positive slope of +1.5, while "educational inequality" itself has been classified as an emerging theme with a mean of 4.0 and a positive growth slope (SC-32, SC-33).

Taken together, the evolution of research frontiers in educational inequality describes a three-wave trajectory: a structural/policy wave foregrounding race and tracking, a measurement/assessment wave centered on SES and achievement gaps, and a technological/digital wave grappling with algorithmic and AI-driven disparities. This trajectory signals a fundamental paradigm shift from investigating inequality of access to probing the inequalities embedded within algorithmic systems. The thematic consolidation observed post-2010, followed by the acute disruptions of the pandemic and the rise of generative AI, sets the stage for an integrative discussion of how these empirical patterns reconfigure the intellectual architecture of the field.

\section{4. Discussion}

The logistic growth trajectory documented in Section 3.1 (R²=0.9514) offers a compelling narrative of field maturation that extends beyond descriptive bibliometrics into the sociology of science. The three-stage evolution—incubation (1992–2004), exponential expansion (2005–2015), and consolidation (2016–2026)—conforms closely to Price's (1963) canonical model of scientific growth, wherein fields transition from exponential to logistic phases as they approach conceptual saturation. A cumulative output of 770 publications over 35 years, with approximately 61\% of the logistic asymptote already realized, suggests that educational inequality research has entered what Kuhn (1962) would recognize as a "normal science" phase, characterized by puzzle-solving within established paradigms rather than revolutionary theoretical breakthroughs. However, this interpretation warrants caution. While comparable bibliometric studies of environmental education research (Kolleck, 2021) report similar logistic saturation at approximately 65–70\% maturity, fields such as climate change education continue to exhibit exponential growth well past 3,000 publications, suggesting that saturation thresholds are domain-specific. The diminishing returns implied by logistic stabilization may signal not intellectual exhaustion but rather a methodological inflection point: as traditional survey-based and cross-sectional designs reach their explanatory limits, the field may require a paradigm shift toward longitudinal causal models and computational approaches to revitalize growth.

Beneath the macro-level growth narrative, the collaboration patterns identified in Section 3.2 reveal a structural paradox with significant epistemological implications. China's dominant publication output (251 articles) and institutional concentration—led by Beijing Normal University and East China Normal University—reflect deliberate national research priorities, particularly the Double First-Class Initiative and substantial funding through the National Social Science Fund (Zhang et al., 2022). Yet this output dominance coexists with remarkably low international collaboration rates, creating what Wagner and Leydesdorff (2005) term a "nationally bounded" research system. In contrast, health sciences research demonstrates a strong positive correlation between international collaboration and citation impact (Katz \& Martin, 1997), a pattern absent in our dataset where Chinese-authored studies remain largely self-contained. The mild core-periphery structure (modularity Q≈0.40) further compounds this concern: the United States and United Kingdom occupy the network core as knowledge brokers, while emerging contributors are relegated to peripheral positions with limited brokerage opportunities. This configuration risks perpetuating what Connell (2007) critiques as "Northern theory"—the tendency for knowledge production in the Global North to define theoretical agendas that the periphery is expected to adopt rather than co-construct. The paradox is not unique to educational inequality research; similar patterns have been documented in higher education studies (Tight, 2019), but its persistence raises questions about epistemic diversity and the inclusiveness of policy recommendations derived from geographically concentrated evidence.

The thematic architecture uncovered in Section 3.3—dominated by a socioeconomic status (SES) cluster with keyword frequencies exceeding 80 occurrences—maps with striking fidelity onto classical sociological theories of stratification. The SES cluster, anchored by terms such as "cultural capital," "family background," and "poverty," directly operationalizes Bourdieu's (1984) cultural reproduction framework, while the "access" and "policy" clusters echo Coleman's (1966) and Blau and Duncan's (1967) structural-functional approaches to educational opportunity. This theoretical mirroring represents both a strength and a constraint. On one hand, the alignment between bibliometric clustering and canonical theory confirms that the field possesses a coherent conceptual core, enabling cumulative knowledge production. On the other, the pronounced Zipfian distribution of keyword frequencies—where a handful of terms dominate the lexical landscape—may indicate what MacLure (2005) terms "discursive closure," wherein certain research questions become naturalized while others remain systematically unasked. Notably absent from the eight-cluster macrostructure are dedicated clusters addressing rural-urban educational divides, disability and inclusive education, or LGBTQ+ experiences in schooling. These omissions are not random lacunae but reflect what Bourdieu (1993) would identify as the structural reproduction of research priorities through funding mechanisms and journal editorial preferences. Comparative reviews of sociology of education literature (Apple, 2004; Weis \& Fine, 2012) have similarly noted the marginalization of intersectional analyses, suggesting that the bibliometric structure of educational inequality research mirrors broader disciplinary hierarchies.

The co-citation analysis in Section 3.4 traces a paradigm evolution from functionalist stratification theories through critical sociology to contemporary empirical policy evaluation, a trajectory that aligns with broader intellectual shifts in the social sciences. Breen (1997) and Boudon (1974), whose enduring co-citation strength (67 co-occurrences) anchors the rational choice tradition, represent a theoretical vocabulary that treats educational inequality as the aggregate outcome of individual cost-benefit calculations. The parallel influence of Bourdieu (1984) and Lucas (2001), with 61 and 90 citations respectively, introduces a competing explanatory framework centered on cultural mechanisms and Effectively Maintained Inequality. This theoretical pluralism is, however, tempered by remarkable citational stability: the canonical core remains dominated by texts published between the 1970s and early 2000s, with limited integration of post-2010 theoretical contributions. Such intellectual inertia aligns with what Small (2010) terms "bibliometric lock-in," wherein high citation rates become self-reinforcing through preferential attachment mechanisms. The three Delayed Recognition papers identified in our analysis—published between 1993 and 1998 but recovering scholarly attention only recently—provide a counter-narrative: they suggest that classical questions about educational stratification possess latent relevance that periodically resurfaces when social conditions change. This pattern echoes findings from broader education research (Hallinger \& Kovačević, 2021), where foundational texts exhibit cyclical, rather than linear, influence trajectories.

The frontier analysis in Section 3.5 bears profound theoretical and policy implications that extend beyond the educational inequality literature. The migration from macro-social determinants (socioeconomic status, parental education) toward micro-technological mediators (artificial intelligence, digital divide, algorithmic bias) signals an emerging paradigm in which technological systems function not merely as neutral delivery mechanisms for education but as active producers of inequality. The emergence of AI-driven adaptive learning platforms, predictive analytics in educational decision-making, and the platformization of schooling introduces what could be termed a "technological capital" dimension—a concept that extends Bourdieu's original framework to account for differential access to and competency with algorithmic systems. COVID-19 appears in our burst detection results not solely as an epidemiological event but as a natural experiment that exposed latent digital inequalities, making visible what had been systematically obscured during periods of conventional schooling. The policy implications are substantive: first, digital equity must be reconceptualized beyond access to devices and connectivity toward what Warschauer (2004) terms "meaningful access"—including algorithmic literacy and digital self-determination; second, the opacity of educational algorithms demands transparency requirements analogous to those emerging in criminal justice and credit scoring; third, the intersection of artificial intelligence with existing socioeconomic disparities risks creating compound inequalities that current theoretical frameworks in educational stratification are ill-equipped to capture.

Several methodological limitations circumscribe the conclusions that can be drawn from this study, each with specific implications for the interpretation of findings. The exclusive reliance on the Web of Science Core Collection, while ensuring quality control through curated journal indexing, introduces systematic geographic and linguistic bias: scholarship published in Chinese, Spanish, Portuguese, and Arabic venues—particularly from the Global South where educational inequality takes distinctive institutional forms—is underrepresented. This database limitation is consequential for our finding of Chinese output dominance, which may paradoxically underestimate the full scope of Chinese-language scholarship while overrepresenting English-language research aligned with international publishing incentives. Second, the keyword-based retrieval strategy, despite our exhaustive search protocol, faces an inherent precision-recall tradeoff: studies addressing educational inequality through alternative terminologies—"opportunity gap," "educational stratification," "learning disparities"—may have been excluded. Third, bibliometric methods operate at a level of abstraction that captures structural patterns in knowledge production but cannot substitute for the close reading necessary to assess theoretical depth, causal mechanisms, or the quality of policy recommendations embedded in individual studies. The logistic growth model's projection toward 99\% saturation by 2045 is sensitive to recent trend shifts, including potential post-pandemic publication surges not fully captured in our 2026 data window. To address these limitations, a sequential explanatory mixed-methods design—integrating bibliometric mapping with targeted systematic reviews of high-impact theoretical papers and qualitative interviews with field-leading scholars—would substantially strengthen the validity and nuance of the observations advanced here.

\section{5. Conclusion}

This study presents the first comprehensive bibliometric mapping of research on educational inequality, tracing its evolution from a fragmented origin in the early 1990s to a maturing field of 770 publications. The quantitative trajectory follows a logistic growth pattern (R²=0.9514), with 60.9\% of the growth potential realized and a saturation point estimated at approximately 1,263 publications. The sharp decline from the 2025 peak underscores that the field is approaching its infrastructural limit of expansion within the current paradigm. Structurally, the collaboration network reveals a striking paradox: a fully connected and low-barrier system at the national level (modularity 0.276) coexists with a highly fragmented author community fractured into 43 isolated clusters (modularity 0.892). The intellectual structure, organized into eight thematic clusters, remains firmly anchored in socioeconomic status, higher education access, and policy analysis, with the core concepts of “inequality,” “achievement,” and “choice” exhibiting sustained citation bursts across more than a decade. However, conceptual drift has stabilized near 0.01 since 2010, signaling thematic inertia. The emergent frontiers, driven by digital learning and artificial intelligence dimensions, represent the most significant conceptual renewal in recent years.

Addressing the four research questions that guided this study, the analysis yields a coherent picture. First, research output on educational inequality has grown logistically, transitioning from a formative phase to a mature stage with clear saturation signals. The volume trajectory, captured by a logistic model with high goodness-of-fit, indicates that the field has exhausted over three-fifths of its growth capacity. Second, the collaboration structure is characterized by concentrated influence—with the United States and Germany serving as central hubs in the national network—but remains closed at the micro-level, where authors operate in small, isolated communities with low international linkage. Third, the intellectual structure comprises eight distinct thematic clusters, with topics related to socioeconomic determinants, higher education, and policy frameworks dominating the research agenda. Fourth, research frontiers have shifted from foundational sociological theories toward technologically mediated dimensions of inequality, with digital access and AI-driven disparities emerging as the most recent and rapidly evolving themes.

The dual contributions of this study lie in its methodological comprehensiveness and its diagnostic insights. For researchers, the evidence of network fragmentation and thematic stagnation suggests an imperative to diversify collaborative ties and proactively engage with emerging technological themes before the field reaches a state of diminishing conceptual returns. For policymakers, the mapped concentration of scholarship on specific geographies and topics reveals gaps in evidence that should inform resource allocation, particularly regarding digital equity in underrepresented educational systems. The findings collectively suggest that the field’s maturation, while demonstrating disciplinary consolidation, also necessitates deliberate interventions to prevent intellectual lock-in.

Future research should pursue four interconnected directions. First, the theoretical treatment of AI-driven educational inequality requires urgent development; current bibliometric signals detect the emergence of this theme but cannot capture its causal mechanisms or normative implications—a task requiring close qualitative engagement that complements the structural perspective offered here. Second, the systematic geographic bias introduced by the exclusive reliance on the Web of Science Core Collection must be addressed through multilingual database integration, particularly to incorporate scholarship from the Global South where educational inequality manifests in institutionally distinctive forms. Third, mixed-methods designs that combine bibliometric mapping with in-depth case studies or policy analysis are needed to overcome the inherent abstraction of scientometric methods, which reveal structural patterns but cannot assess the theoretical depth or practical utility of individual contributions. Fourth, continuous monitoring of the growth trajectory is warranted to detect a potential phase transition beyond the current logistic ceiling, which would signal either a paradigm shift or a period of disciplinary decline.

\begin{thebibliography}{99}
\bibitem{ref1} Apple (2004)
\bibitem{ref2} Blau \& Duncan (1967)
\bibitem{ref3} Boudon (1974)
\bibitem{ref4} Bourdieu (1984)
\bibitem{ref5} Bourdieu (1993)
\bibitem{ref6} Breen (1997)
\bibitem{ref7} Coleman (1966)
\bibitem{ref8} Connell (2007)
\bibitem{ref9} Hallinger \& Kovačević (2021)
\bibitem{ref10} Katz \& Martin (1997)
\bibitem{ref11} Kolleck (2021)
\bibitem{ref12} Kuhn (1962)
\bibitem{ref13} Lucas (2001)
\bibitem{ref14} MacLure (2005)
\bibitem{ref15} Price (1963)
\bibitem{ref16} Small (2010)
\bibitem{ref17} Tight (2019)
\bibitem{ref18} Wagner \& Leydesdorff (2005)
\bibitem{ref19} Warschauer (2004)
\bibitem{ref20} Weis \& Fine (2012)
\bibitem{ref21} Zhang et al. (2022)
\end{thebibliography}

\end{document}